

\newcommand{\comment}[1]{}
\newcommand{\myBigBreak}{\\[5mm]}


\label{01-chapter}

%%%%% Capitolo 2 %%%%%

\section{Introduzione al capitolo} \label{sezione 1-1}
TODO descrizione del capitolo\par





\section{Che cos'è il Machine Learning} \label{sezione 1-2}
\comment{
Imagine you own a machine that manufactures widgets. Sometimes it breaks down,
and it’s expensive to repair. Perhaps if you collected data about the machine during
operation, you might be able to predict when it is about to break down and halt oper‐
ation before damage occurs. For instance, you could record its rate of production, its
temperature, and how much it is vibrating. It might be that some combination of
these factors indicates an impending problem. But how do you figure it out?
This is an example of the sort of problem machine learning is designed to solve. Fun‐
damentally, machine learning is a technique for using computers to predict things
based on past observations. We collect data about our factory machine’s performance
and then create a computer program that analyzes that data and uses it to predict
future states.
}
/comment{
Prendiamo ad esempio un qualsiasi gesto compiuto con il movimento di una mano, se collezioniamo abbastanza dati durante il movimento, forse è possibile far capire ad una macchina quale movimento è stato realizzato, \comment{ovvero classificare i dati raccolti con uno specifico label, ad esempio tenendo il pollice alzato e il resto della mano chiusa a pugno potremmo ricevere un label "OK", oppure tenendo la mano aperta e scuotendola ricevere "CIAO",} potremmo registrare ad esempio la velocità e accelerazione di alcune dita e del polso, e una combinazione di questi fattori potrebbe portarci alla soluzione che desideriamo, ma in che modo?\\
Questo è un esempio del genere di problemi, che il machine learning è stato progettato per risolvere.\\
Utilizzando le osservazioni passate, fornendo ingressi ed uscite, (riprendendo l'esempio di prima i dati degli accelerometri come ingressi, ed il nome del gesto come uscita) il programma di machine learning è in grado analizzare questi dati, e creare un modello che possa generalizzare a quali uscite corrispondono nuovi ingressi. 
\myBigBreak
}
\comment{
Creating a machine learning program is different from the usual process of writing
code. In a traditional piece of software, a programmer designs an algorithm that takes
an input, applies various rules, and returns an output. The algorithm’s internal opera‐
tions are planned out by the programmer and implemented explicitly through lines
of code. To predict breakdowns in a factory machine, the programmer would need to
understand which measurements in the data indicate a problem and write code that
deliberately checks for them.
}
/comment{
Scrivere un programma di machine learning è molto diverso da scrivere un tipico codice, tradizionalmente dati degli ingressi , otteniamo delle uscite attraverso una serie di funzioni e regole date dal programmatore, i dati in ingresso vengono processati e restituiti come nuovi dati in output. Tutte le funzioni o regole, che decidono come i dati dovranno essere modificate, sono definite dal programmatore, ed implementate esplicitamente attraverso righe di codice.\\
Per capire a quale gesto, corrisponde un vettore di dati ottenuti da vari accelerometri, il programmatore deve capire quali misure corrispondono a quali gesti, e scrivere codice che deliberatamente cerca per queste misure.
\myBigBreak
}

\comment{
This approach works fine for many problems. For example, we know that water boils
at 100°C at sea level, so it’s easy to write a program that can predict whether water is
boiling based on its current temperature and altitude. But in many cases, it can be
difficult to know the exact combination of factors that predicts a given state. To con‐
tinue with our factory machine example, there might be various different combina‐
tions of production rate, temperature, and vibration level that might indicate a
problem but are not immediately obvious from looking at the data.
}

\comment{
Il tradizionale procedimento funziona bene per molti problemi come ad esempio capire se un oggetto è in movimento o meno, semplicemente se l'accelerazione letta dell'accelerometro cambia possiamo definire l'oggetto come in "movimento". Ma in molti altri casi è difficile conoscere quale combinazione di fattori per predire lo stato del sistema, ad esempio quando i dati ricevuti dai sensori sono sull'ordine delle decine, ovvero quando non è possibile facilmente determinare la soluzione guardando i dati.
\myBigBreak
}

\comment{
In caso si scelga di utilizzare un algoritmo di machine learning, non è più il programmatore a creare le regole per il processo dei dati, ma la macchina decide da sola un modello che meglio approssima la relazione tra ingressi ed uscite fornitegli.\\
Il processo di apprendimento, dove l'algoritmo decide i pesi del modello, è appunto la parte di "learning".\\
\myBigBreak
}

\comment{
12 | Chapter 3: Getting Up to Speed on Machine Learning
To create a machine learning program, a programmer feeds data into a special kind of
algorithm and lets the algorithm discover the rules. This means that as programmers,
we can create programs that make predictions based on complex data without having
to understand all of the complexity ourselves. The machine learning algorithm builds
a model of the system based on the data we provide, through a process we call train‐
ing. The model is a type of computer program. We run data through this model to
make predictions, in a process called inference.
There are many different approaches to machine learning. One of the most popular is
deep learning, which is based on a simplified idea of how the human brain might
work. In deep learning, a network of simulated neurons (represented by arrays of
numbers) is trained to model the relationships between various inputs and outputs.
Different architectures, or arrangements of simulated neurons, are useful for different
tasks. For instance, some architectures excel at extracting meaning from image data,
while other architectures work best for predicting the next value in a sequence.
The examples in this book focus on deep learning, since it’s a flexible and powerful
tool for solving the types of problems that are well suited to microcontrollers. It
might be surprising to discover that deep learning can work even on devices with
limited memory and processing power. In fact, over the course of this book, you’ll
learn how to create deep learning models that do some really amazing things but that
still fit within the constraints of tiny devices.
The next section explains the basic workflow for creating and using a deep learning
model.
}






